\clearpage
\SourceLeaf{156}
\begin{center}
[V.]

{\large [KLEINERE BEITRÄGE\\
ZUR METHODE DER KLEINSTEN QUADRATE.]}
\end{center}



$A,\ B,\ B',\ \ldots\ B^{(n-1)}$ seien linearische Functionen von $x$ und
andern unbekannten Grössen.
\[
\begin{array}{rl}
A&=\alpha x+\cdots\\
B&=\beta x+\cdots\\
B'&=\beta' x+\cdots\\
.&\hphantom{=} .\\[-.6em]
.&\hphantom{=} .\\
B^{(n-1)}&=\beta^{(n-1)}x+\cdots.
\end{array}
\]
\[
\Omega = (\lambda A)^2+(\mu B)^2+(\mu' B')^2+(\mu'' B'')^2+\cdots
\]
soll ein Minimum werden.

Man setze
\[
\frac{\alpha}{\beta}=l,\qquad \frac{\mu}{\lambda l}=\mathrm{tang}\,\varphi,\qquad
\frac{\mu\,\mathrm{tang}\,\varphi}{\lambda}=k,
\]
\[
A+kB=A',\qquad A-lB=C,
\]
\[
\lambda\cos\varphi=\lambda',\qquad \lambda\sin\varphi=\nu,
\]
so ist
\[
(\lambda' A')^2+(\nu C)^2=(\lambda A)^2+(\mu B)^2,
\]
und $C$ von $x$ frei.

Der Coefficient von $x$ in $A'$ wird $=\dfrac{\alpha}{\cos\varphi^2}$.

